\paragraph{P1.7：同预算条件 law 与强控制对照。}
复用 P1.6 已计费校准和同一模型训练池；逐模式选择允许的有限映射，
完整方法与 task-only 使用完全相同的容量。额外模型验证单列计数，目标
任务轨迹不参与选择。对六个新 reset groups 使用成对的 antithetic
uniforms，每组源模式恰为两支，源 law 的频率在该条件 fixture 内平衡。
这是 clustered 条件 law 检查，不是独立同分布模式频率估计。
% Generated from independently audited scientific_report.json.
\begin{center}\small
\begin{tabular}{lrr}
\toprule
Method & Derated & Matched\\
\midrule
Full & 28/36 & 23/36\\
Task-only & 28/36 & 23/36\\
Nominal & 8/36 & 23/36\\
Reweight only & 0/36 & 0/36\\
Task-only + reweight & 35/36 & 35/36\\
Shuffled coupling & 28/36 & 23/36\\
Extra noise & 28/36 & 24/36\\
No source mode & 20/36 & 34/36\\
\bottomrule
\end{tabular}
\end{center}
本阶段新探针 0 次，实际不同目标评估 132 条；别名方法与拒绝另报。


v1 full/task-only 六次选择相同条件映射，均 $51/72$，其中 $18$ 个请求
因不支持而拒绝。task-only 加重加权为 $70/72$，说明限制模式质量会影响
任务分母。shuffle 保留同样 $51/72$ 和成功模式数，却没有原始源对应；
额外噪声 $52/72$ 不能作为其必要性的证据。

恢复支持 v2 继续使用同一候选/校准缓冲，并对 full 与 task-only 同时开放
重加权。完整结果如下；不能把 v1/v2 中最好的格拼成正式主表。
% Generated from independently audited scientific_report.json.
\begin{center}\small
\begin{tabular}{lrr}
\toprule
Method & Derated & Matched\\
\midrule
Full & 30/36 & 34/36\\
Task-only & 30/36 & 35/36\\
Nominal & 11/36 & 34/36\\
Reweight only & 0/36 & 0/36\\
Task-only + reweight & 36/36 & 35/36\\
Full + reweight & 36/36 & 34/36\\
Shuffled coupling & 30/36 & 34/36\\
Extra noise & 30/36 & 35/36\\
No source mode & 36/36 & 34/36\\
\bottomrule
\end{tabular}
\end{center}
本阶段新探针 0 次，实际不同目标评估 186 条；别名方法与拒绝另报。

v2 full/task-only 分别为 $64/72$ 与 $65/72$，两者各含六个拒绝；
加重加权后为 $70/72$ 与 $71/72$。matched 的一个新 reset 中，full
在 mode0 丢失支撑而 task-only 成功；共享 mode1 均超力。
该校准重复的实际样本均值 task 差为 $+0.06788$，响应差为 $-0.00309$。
这说明训练预算带并未认证新 reset 的非劣，响应均值下降也不等于任务收益。
当前数据没有建立 full 相对同类控制优化的独立优势。

去掉 source mode 与打乱 coupling 检验源条件对应；它们不能单独证明
response objective 相对同 prior 的 task-only 有独立价值。仅重加权保持
固定命令支持，另加入 task-only 加重加权的强对照。额外 transport noise
固定幅度、保留全部约束，确定性映射仍从随机源 prior 获得多模态。
